\documentclass[11pt,letterpaper]{article}
\def\SuiteDocTitle{Chapter-to-Operations Crosswalk}
\def\SuiteDocCode{CU-02}
% Shared professional style: pdfLaTeX compatible, no external font files.
\usepackage[T1]{fontenc}
\usepackage[utf8]{inputenc}
\usepackage{lmodern}
\usepackage{microtype}
\usepackage[margin=0.83in,headheight=16pt,footskip=28pt]{geometry}
\usepackage[table]{xcolor}
\definecolor{navy}{HTML}{12304A}
\definecolor{teal}{HTML}{087F83}
\definecolor{gold}{HTML}{B78633}
\definecolor{blueback}{HTML}{F0F5F8}
\definecolor{tealback}{HTML}{EEF8F6}
\definecolor{amberback}{HTML}{FFF7E8}
\definecolor{grayback}{HTML}{F6F7F8}
\usepackage{array,tabularx,longtable,booktabs}
\usepackage{enumitem}
\usepackage{amssymb}
\usepackage{titlesec}
\usepackage{fancyhdr,lastpage}
\usepackage{tcolorbox}
\tcbuselibrary{breakable}
\usepackage{pgffor}
\usepackage{hyperref}
\hypersetup{colorlinks=true,linkcolor=teal,urlcolor=teal,pdfauthor={Engineering Learning Series},pdftitle={\SuiteDocTitle}}
\setlength{\parindent}{0pt}
\setlength{\parskip}{5.5pt}
\setlist{topsep=3pt}
\setcounter{tocdepth}{2}
\titleformat{\section}{\Large\bfseries\color{navy}}{\thesection}{0.7em}{}
\titleformat{\subsection}{\large\bfseries\color{teal}}{\thesubsection}{0.6em}{}
\titleformat{\subsubsection}{\normalsize\bfseries\color{navy}}{\thesubsubsection}{0.6em}{}
\pagestyle{fancy}
\fancyhf{}
\fancyhead[L]{\scriptsize\color{navy}AGILE REQUIREMENTS ENGINEERING}
\fancyhead[R]{\scriptsize\color{teal}\SuiteDocCode}
\fancyfoot[L]{\scriptsize Independent learning and implementation materials}
\fancyfoot[R]{\scriptsize Page \thepage\ of \pageref{LastPage}}
\renewcommand{\headrulewidth}{0.35pt}
\renewcommand{\footrulewidth}{0pt}
\newcommand{\SuiteHeader}[2]{%
{\color{teal}\small\bfseries ENGINEERING DOCUMENTATION SUITE \hfill #2}\par\vspace{8pt}
{\color{navy}\LARGE\bfseries\SuiteDocTitle}\par\vspace{5pt}
{\color{black!68}\normalsize #1}\par\vspace{8pt}
{\color{gold}\rule{\linewidth}{2pt}}\par\vspace{8pt}}
\newcommand{\WritingLines}[1]{\par\foreach \n in {1,...,#1}{\noindent\rule{\linewidth}{0.35pt}\par\vspace{0.28cm}}}
\newcommand{\FreeSpace}[1]{\begin{tcolorbox}[colback=white,colframe=black!25,height=#1,boxrule=0.5pt]\end{tcolorbox}}
\newcommand{\Field}[1]{\noindent\textbf{#1}\quad\rule{0.42\linewidth}{0.35pt}\par\vspace{0.3cm}}
\newcommand{\CheckItem}[1]{\noindent$\square$\quad #1\par\vspace{2pt}}
\newcommand{\SmallNote}[1]{{\footnotesize\color{black!62}#1}\par}
\newcommand{\Guidance}[1]{\begin{tcolorbox}[colback=blueback,colframe=teal,breakable,title=Guidance]#1\end{tcolorbox}}
\sloppy

\begin{document}
\SuiteHeader{Book chapter index | workflows, artifacts and exercises}{Reference}
\tableofcontents\medskip\noindent\rule{\linewidth}{0.4pt}\medskip
\section{How to use the crosswalk}
Printed book page numbers refer to the 2004 Addison-Wesley edition in the supplied PDF; PDF viewer page indices differ because of front matter. Each chapter maps to this suite’s relevant exercises and original implementation artifacts. The crosswalk is a navigation aid, not a reproduced table of contents.

\begin{longtable}{|p{0.25\linewidth}|p{0.35\linewidth}|p{0.22\linewidth}|}\hline
\textbf{Chapter / printed pages} & \textbf{Learning focus} & \textbf{Suite application} \\ \hline\endfirsthead\hline
\textbf{Chapter / printed pages} & \textbf{Learning focus} & \textbf{Suite application} \\ \hline\endhead
1. An Overview (3–15) & Three-part story discipline; Shared decisions under uncertainty; Planning at multiple horizons & Study guide SG-01 \\ \hline
2. Writing Stories (17–29) & INVEST as a review heuristic; Value belongs to an identifiable role; Size and testability & Study guide SG-02 \\ \hline
3. User Role Modeling (31–41) & Role not generic user; Role consolidation and modeling; Personas as optional enrichment & Study guide SG-03 \\ \hline
4. Gathering Stories (43–53) & Elicitation through multiple channels; Goal-oriented questions; Just-enough capture & Study guide SG-04 \\ \hline
5. Working with User Proxies (55–66) & Proxy varieties; Proxy bias; Customer team construction & Study guide SG-05 \\ \hline
6. Acceptance Testing User Stories (67–74) & Confirm intent before coding; Customer intent and engineering collaboration; Right-sized test selection & Study guide SG-06 \\ \hline
7. Guidelines for Good Stories (75–83) & Start with goals; split vertically; Closed stories and constraints; Defer premature design detail & Study guide SG-07 \\ \hline
8. Estimating User Stories (87–95) & Relative story points; Whole-team estimation; Triangulation & Study guide SG-08 \\ \hline
9. Planning a Release (97–107) & Release tradeoffs; Risk and infrastructure; Velocity-informed projection & Study guide SG-09 \\ \hline
10. Planning an Iteration (109–116) & Reopen story conversations; Story-to-task decomposition; Capacity and ownership & Study guide SG-10 \\ \hline
11. Measuring and Monitoring Velocity (117–129) & Complete stories only; Burndown as a signal; Progress conversations & Study guide SG-11 \\ \hline
12. What Stories Are Not (133–143) & Stories versus detailed specifications; Stories versus use cases; Stories versus scenarios & Study guide SG-12 \\ \hline
13. Why User Stories? (145–155) & Conversation and comprehensibility; Late detail, early value; Limits of the approach & Study guide SG-13 \\ \hline
14. A Catalog of Story Smells (157–163) & Granularity smells; Excess and premature detail; Weak product participation & Study guide SG-14 \\ \hline
15. Using Stories with Scrum (165–176) & Product backlog alignment; Sprint goal and inspection; Cross-functional ownership & Study guide SG-15 \\ \hline
16. Additional Topics (177–187) & Nonfunctional constraints; Interface work and artifact persistence; Defects and emergent work & Study guide SG-16 \\ \hline
17. The User Roles (191–198) & Example context and stakeholder discovery; Role narrowing; Persona linkage & Study guide SG-17 \\ \hline
18. The Stories (199–206) & Role-specific story inventory; Discovery before full specification; Coverage across goals & Study guide SG-18 \\ \hline
19. Estimating the Stories (209–218) & Example-driven relative estimation; Uncertainty in estimates; Backlog estimate consistency & Study guide SG-19 \\ \hline
20. The Release Plan (219–222) & Estimate capacity and select outcomes; Manage risk visibly; Adapt as evidence changes & Study guide SG-20 \\ \hline
21. The Acceptance Tests (223–229) & Case-study confirmation; Traceability; Release evidence & Study guide SG-21 \\ \hline
\end{longtable}
\section{Book-to-operation crosswalk}
\begin{longtable}{|p{0.20\linewidth}|p{0.29\linewidth}|p{0.37\linewidth}|}\hline
\textbf{Operational activity} & \textbf{Core chapters} & \textbf{Primary suite artifacts} \\ \hline\endfirsthead\hline
\textbf{Operational activity} & \textbf{Core chapters} & \textbf{Primary suite artifacts} \\ \hline\endhead
Discovery and role modeling & 1, 3, 4, 5, 17 & Role model, discovery workbook, Lab 01 \\ \hline
Authoring and refinement & 2, 7, 14, 18 & Story card, slicing workbook, INVEST scorecard, Labs 02–04 \\ \hline
Acceptance and quality & 6, 16, 21 & Testing handbook, acceptance workbook, Lab 05 \\ \hline
Estimation and forecasting & 8, 9, 10, 11, 19, 20 & Release workbook, planning templates, Labs 06–07 \\ \hline
Scrum implementation & 10, 15 & Iteration planning procedure and template \\ \hline
Formal requirements and governance & 12, 13, 16 & Traceability handbook, change record, Lab 08 \\ \hline
Integrated delivery & 17–21 & Lab 09 capstone and scorecard package \\ \hline
\end{longtable}
\section{Recommended progression}
First read 1–7 to establish discovery and confirmation. Next read 8–11 for estimation and planning. Study 12–16 to judge boundaries, smells, Scrum and constraints. Use 17–21 as a continuous worked example, returning to earlier chapters whenever a decision is unclear.


\end{document}
